% Cumplimiento del Cronograma (Actividades 1--8)
\chapter{Cumplimiento del Cronograma}
\label{chap:cumplimiento-cronograma}

Este capítulo resume \textit{cómo} se ejecutó cada actividad del cronograma y \textit{por qué} se logró satisfactoriamente, con referencias a la evidencia en el repositorio.

\section*{Resumen}
Todas las actividades (1--8) se marcaron como COMPLETADAS (ver `docs/Macroproceso.txt`, Registro de cumplimiento).

\section{Actividad 1: Revisión de Literatura y Definición de Metodología}
\textbf{Cómo}: Revisión sistemática y análisis documental; definición de enfoque mixto, diseño (experimental/cuasi) y operacionalización (VI/VD, energía J/MB). Redacción de metodología.\\
\textbf{Por qué}: Objetivos claros y alineados al Macroproceso; especificaciones internas de monitoreo/energía/arquitectura.\\
\textbf{Evidencia}: `docs/achievements/metodologia.{md,tex}`, `docs/monitoring_spec.md`, `docs/energy_plan.md`, `docs/architecture/README.md`.

\section{Actividad 2: Diseño de la Arquitectura del SMA y del Experimento}
\textbf{Cómo}: Modelado de agentes e interacciones; puente JADE; flujo PREPROCESS$\to$TRAIN$\to$EVAL; observabilidad con Prometheus/Grafana; Compose y smoke tests.\\
\textbf{Por qué}: Modularidad y separación de preocupaciones; infraestructura y dashboards versionados.\\
\textbf{Evidencia}: `docs/architecture/*`, `infra/compose/*`, `docs/grafana/README.md`, `docs/grafana/thesis_minimal.json`.

\section{Actividad 3: Implementación y Desarrollo de Prototipos}
\textbf{Cómo}: Orquestador JADE; backend FastAPI con métricas; entrenamiento RF/SVM; scripts de ejecución y simulación.\\
\textbf{Por qué}: Integración clara entre agentes, backend y telemetría; toolchain reproducible.\\
\textbf{Evidencia}: `agents/jade-platform/*`, `backend/api/main.py`, `python-analysis/simulate_experiment.py`, `scripts/run_experiment.*`, `scripts/run_batch.*`, `scripts/monitoring_smoke.sh`.

\section{Actividad 4: Recolección y Pre-procesamiento de Datos Experimentales}
\textbf{Cómo}: Dataset sintético reproducible; corridas control/tratamiento; registro de meta y eventos; consolidación en CSV.\\
\textbf{Por qué}: Procedimientos estandarizados y estructura de resultados por ejecución/lote.\\
\textbf{Evidencia}: `python-analysis/generate_dataset.py`, `data/raw/*`, `data/results/<exp_id>/*`, `python-analysis/aggregate_metrics.py`, `data/results/aggregate_metrics.csv`.

\section{Actividad 5: Análisis e Interpretación de Datos}
\textbf{Cómo}: Agregación y análisis estadístico (descriptivo e inferencial); reportes HTML y resúmenes.\\
\textbf{Por qué}: Pipeline analítico versionado y reproducible; métricas homogéneas y normalizadas.\\
\textbf{Evidencia}: `python-analysis/aggregate_metrics.py`, `python-analysis/stats_analysis.py`, `data/results/aggregate_metrics.csv`, `data/results/stats_summary.{json,md}`, `python-analysis/report_run.py`.

\section{Actividad 6: Visualización de Hallazgos y Elaboración de Modelos}
\textbf{Cómo}: Dashboards Grafana (CPU/Mem/Throughput/Energía/progreso); guía de uso; figuras para tesis.\\
\textbf{Por qué}: Dashboards/datasource versionados y métrica en tiempo real.\\
\textbf{Evidencia}: `docs/grafana/README.md`, `docs/grafana/thesis_minimal.json`, `docs/grafana/img/*`, `python-analysis/generate_thesis_figures.py`.

\section{Actividad 7: Redacción del Reporte de Resultados}
\textbf{Cómo}: Redacción de capítulos (MD + LaTeX) con evidencias, incidencias y validez; integración de resultados y gráficos.\\
\textbf{Por qué}: Base analítica consolidada; flujo reproducible y activos versionados.\\
\textbf{Evidencia}: `docs/thesis_results_section.{md,tex}`, `docs/achievements/cumplimiento_objetivo_general.{md,tex}`, `docs/achievements/objetivos_especificos.{md,tex}`.

\section{Actividad 8: Consolidación de Conclusiones, Implicaciones y Trabajos Futuros}
\textbf{Cómo}: Síntesis de hallazgos; discusión de limitaciones, amenazas a la validez e implicaciones; roadmap.\\
\textbf{Por qué}: Evidencia cuantitativa/cualitativa suficiente y documentación continua.\\
\textbf{Evidencia}: `docs/conclusions_implications_future_work.md`, `docs/future_work_roadmap.md`, `docs/resilience_test_scenarios.md`.

\section*{Factores transversales de éxito}
Planificación y trazabilidad; reproducibilidad (scripts/datos); observabilidad (Prometheus/Grafana); normalizaciones energéticas (watts-first, J/MB).
